\section{Algorithm: Formula-Guided Onset Search}
\label{sec:algorithm}

The analysis of \S\ref{sec:analysis} reduces optimisation to \emph{onset detection}:
because $f(M)$ rises to a plateau (\textbf{P1}), the operating point of interest is the
smallest $M$ that reaches near-peak throughput,
\begin{equation}
  \hat M \;=\; \min\{\, M : f(M) \ge (1-\varepsilon)\,f^\ast \,\},
  \qquad \varepsilon = 0.02,
  \label{eq:objective}
\end{equation}
where the peak $f^\ast$ is not known in advance and must itself be probed. Two naive
strategies bracket the design space. Returning $M_{\max}$ costs a single oracle call but
over-provisions concurrency, defeating the purpose of optimising the onset. A
parameter-blind search that simply \emph{maximises} $f$ converges to a point on the
plateau interior --- the wrong objective for Eq.~\eqref{eq:objective} --- and spends many
oracle calls doing so. We instead exploit the closed-form structure of
\S\ref{sec:analysis}: warm-start from the closed-form primitives (\textbf{P2}), then run a
short downward monotone-predicate search (\textbf{P3}) whose threshold is anchored to avoid
the occupancy-gap trap (\textbf{P4}). Algorithm~\ref{alg:onset} states the procedure.

\begin{algorithm}[t]
  \DontPrintSemicolon
  \SetKwFunction{FBisect}{Bisect}
  \SetKwProg{Fn}{Function}{:}{}
  \KwIn{params; oracle $f(\cdot)$; bounds $m_{\min}, m_{\max}$; tolerance $\varepsilon$}
  \KwOut{onset estimate $\hat M \in [m_{\min}, m_{\max}]$}
  $M_{\mathrm{ITL}} \leftarrow$ largest $B$ with $\mathrm{itl}(B) \le \mathrm{targetITL}$
    \tcp*{closed form, no oracle calls}
  $M_{\mathrm{TPF}} \leftarrow$ largest $B$ with $\mathrm{ttft\_prefill}(B) \le \mathrm{targetTTFT}$\;
  \eIf{$M_{\mathrm{TPF}} < M_{\mathrm{ITL}}$}{
    $\mathit{seed} \leftarrow m_{\max}$ \tcp*{ttft-only cell: $M_{\mathrm{TPF}}$ undershoots}
  }{
    $\mathit{seed} \leftarrow \min(\max(M_{\mathrm{ITL}}, M_{\mathrm{TPF}}), m_{\max})$\;
  }
  $f^\ast \leftarrow f(\mathit{seed})$ \tcp*{1 oracle call}
  \If{$f^\ast \le 0$ \textbf{and} $\mathit{seed} < m_{\max}$}{
    $\mathit{seed} \leftarrow m_{\max}$;\quad $f^\ast \leftarrow f(\mathit{seed})$\;
  }
  \lIf{$f^\ast \le 0$}{\Return $m_{\max}$ \tcp*[f]{infeasible everywhere probed}}
  $\theta \leftarrow (1 - \varepsilon/2)\, f^\ast$ \tcp*{margin keeps a confirmatory call $<\varepsilon$}
  $\mathit{lo} \leftarrow \max(m_{\min},\, \min(M_{\mathrm{ITL}}, M_{\mathrm{TPF}}))$\;
  \eIf{$M_{\mathrm{TPF}} < M_{\mathrm{ITL}}$}{
    $\mathit{hi} \leftarrow \min(\mathit{seed}, m_{\max})$\;
  }{
    $\mathit{hi} \leftarrow \min(\mathit{seed},\, 3 M_{\mathrm{ITL}},\, m_{\max})$
      \tcp*{occupancy-gap bound}
  }
  $\mathit{hi} \leftarrow \max(\mathit{lo}, \mathit{hi})$\;
  \Return \FBisect{$f$, $\theta$, $\mathit{lo}$, $\mathit{hi}$}
    \tcp*{$\mathit{lo},\mathit{hi}\in[m_{\min},m_{\max}]$ by construction}
  \BlankLine
  \Fn{\FBisect{$f$, $\theta$, $\mathit{lo}$, $\mathit{hi}$}}{
    \tcp{smallest $M\in[\mathit{lo},\mathit{hi}]$ with $f(M)\ge\theta$, in $\le\textsc{MaxIters}$ probes}
    $i \leftarrow 0$\;
    \While{$\mathit{lo} < \mathit{hi}$ \textbf{and} $i < \textsc{MaxIters}$}{
      $\mathit{mid} \leftarrow \lfloor (\mathit{lo}+\mathit{hi})/2 \rfloor$\;
      \eIf{$f(\mathit{mid}) \ge \theta$}{$\mathit{hi} \leftarrow \mathit{mid}$}{$\mathit{lo} \leftarrow \mathit{mid}+1$}
      $i \leftarrow i + 1$\;
    }
    \Return $\mathit{hi}$\;
  }
  \caption{\textsc{FormulaGuidedOnsetSearch} ($\textsc{MaxIters}=6$); the
           bracketed search is \textsc{Bisect}, a bounded monotone-predicate
           binary search.}
  \label{alg:onset}
\end{algorithm}

\subsection{Why each step}
\label{sec:algo-rationale}

\paragraph{Closed-form classification (P2).}
The primitives $M_{\mathrm{ITL}}$ and $M_{\mathrm{TPF}}$ are computed from the service
model alone, at zero oracle cost (\S\ref{sec:mitl}, \S\ref{sec:regimes}), and their
\emph{ordering} sorts the scenario into a regime cell: TTFT-only when
$M_{\mathrm{TPF}} < M_{\mathrm{ITL}}$, ITL-or-crossover otherwise. This classification reads
only workload and server parameters --- the SLO targets, the service coefficients, and the
batch-size ceiling $m_{\max}$ --- and none of the search constants ($\varepsilon$,
$\textsc{MaxIters}$, the bracket factor). The algorithm's discretion enters only after the
cell is fixed, so the regime structure is a property of the deployment, not an artefact of
the solver.

\paragraph{Cell-dependent seed (P4).}
The seed must read a \emph{near-peak} $f^\ast$: if it underestimates the peak, the threshold
$\theta$ is set too low and the downward search halts prematurely. Which $M$ delivers a
near-peak reading depends on the cell. In the ITL-or-crossover cell the binding ITL
constraint has already plateaued by the upper endpoint
$U=\max(M_{\mathrm{ITL}},M_{\mathrm{TPF}})$, so $f(U)\approx f^\ast$ and we seed at $U$. In
the TTFT-only cell the closed-form endpoints \emph{undershoot}: the prefill bound
$M_{\mathrm{TPF}}$ excludes the M/M/1 queue wait, so the realised TTFT frontier lies above
it (\S\ref{sec:regimes}), and in the worst case $U$ collapses far below the plateau --- on
one benchmark scenario $U = 9$ while peak throughput is attained at $M \approx 57$,
$6.3\times$ higher. Seeding at $U$ there would set $\theta$ off $f(U)\ll f^\ast$ and stop
the descent at $M = 9$, a $22\%$ shortfall that violates the $\varepsilon = 2\%$ tolerance.
We therefore seed at $m_{\max}$ in the TTFT-only cell, which always lies on the plateau; a
final guard re-probes at $m_{\max}$ if even the cell-chosen seed is infeasible.

\paragraph{Bracket.}
The lower end $\mathit{lo} = \min(M_{\mathrm{ITL}}, M_{\mathrm{TPF}})$ floors the search at
$M_{\mathrm{TPF}}$ on the TTFT-only cell, where $M_{\mathrm{ITL}}$ overshoots the onset.
The upper end is tightened by the occupancy-gap bound $3M_{\mathrm{ITL}}$ (the gap stays
within this factor on the scenarios studied), relaxed to the seed on the TTFT-only cell
where $M_{\mathrm{ITL}}$ is an unreliable tightener.

\paragraph{Threshold margin.}
Searching to $\theta = (1-\varepsilon/2)f^\ast$ rather than $(1-\varepsilon)f^\ast$ leaves
headroom so that a downstream confirmatory evaluation of the returned $\hat M$ still lands
within the $\varepsilon$ SLO-feasibility tolerance.

Figure~\ref{fig:onset-search} illustrates the procedure on the baseline scenario (an
ITL-or-crossover cell, so $f^\ast$ is read at $U = M_{\mathrm{TPF}} = 147$): the closed-form
markers, the anchored threshold, and the downward probes descending from the bracket
$[40, 120]$ toward the onset $M^\ast = 69$.

\begin{figure}[t]
  \centering
  \includegraphics[width=0.85\linewidth]{fig5_onset_search}
  \caption{Formula-guided onset search on the baseline (ITL-or-crossover cell). Closed-form
           $M_{\mathrm{ITL}}$ and $M_{\mathrm{TPF}}$ bracket the search; $f^\ast$ is read at
           the constraint endpoint $U = M_{\mathrm{TPF}} = 147$, which already lies on the
           plateau, and the threshold is set at $(1-\varepsilon/2)f^\ast$; the downward
           monotone-predicate probes descend toward the onset $M^\ast$. Illustrative; the
           measured evaluation is in \S\ref{sec:experiments}.}
  \label{fig:onset-search}
\end{figure}

\subsection{Guarantees}
\label{sec:algo-guarantees}

\paragraph{Cost.}
The procedure issues at most two seed probes --- the cell-chosen seed, plus an $m_{\max}$
fallback when that seed is infeasible --- then at most $\textsc{MaxIters}=6$ search probes
and one downstream confirmatory call: nine oracle calls in the worst case, and eight
whenever the seed is feasible, as on every benchmark scenario (RP-8). The search is
$O(\log m_{\max})$ (\textbf{P3}).

\paragraph{Correctness.}
Because $f$ is non-decreasing to a plateau (\textbf{P3}), \textsc{Bisect} returns the
smallest $M$ in the bracket with $f(M) \ge \theta$, so the result meets the near-peak
objective of Eq.~\eqref{eq:objective}. The single high anchor is safe while
the post-peak decline of $f$ stays below $\varepsilon/2$ (\S\ref{sec:search}); under
chunk-count jumps that produce a larger decline, a multi-probe $f^\ast$ would be required
(left to future work).

\paragraph{Onset, not argmax.}
The search deliberately targets the onset $M^\ast$, not the strict throughput argmax, which
can sit far higher on a plateau that wiggles at floating-point precision
(\S\ref{sec:search}). Over-provisioning to the argmax wastes concurrency; any residual gap
between the returned onset and the strict argmax is a structural property of the plateau,
not a defect of the search.

\subsection{Baselines}
\label{sec:algo-baselines}

We compare against two strategies that span the cost/quality extremes. \emph{Naive-max}
returns $M_{\max}$ in a single oracle call but over-provisions concurrency.
\emph{Parameter-blind ternary search} maximises $f$ over $[m_{\min}, m_{\max}]$; on a
monotone-to-plateau curve it converges to an interior plateau point rather than the onset,
solving the wrong objective, and spends roughly thirty oracle calls doing so. The
formula-guided search Pareto-dominates both on the concurrency objective; the measured
comparison across the benchmark is reported in \S\ref{sec:experiments}.

\subsection{Provenance}
\label{sec:algo-provenance}

The closed-form primitives, the three-cell regime partition, and the high-anchor fix were
not hand-designed but \emph{discovered} by an automated formula-discovery campaign that
iteratively proposed and oracle-validated candidate structure against the queueing model.
The campaign's recorded principles (RP-1 through RP-11) are the evidence base for the
design choices above: the exact closed form under $nc=1$ (RP-1), the occupancy-gap
anchor fix (RP-2), the monotone-to-plateau shape (RP-3), the regime partition and its
load-bearing ordering signal (RP-6, RP-7), and the call-budget result (RP-8).
